% % ######% ######% ######% ######% ######% ######% ######% ######% ######% ######
\section{Combined Impact of Work Dealing and Work Stealing}
In this experiment, we evaluate all combinations of work dealing and work stealing together. We want to answer the research question

\begin{tcolorbox}[colback=blue!5, colframe=blue!40!black, title=Research Questions]
    \textbf{RQ:} What is the best combination of work dealing and work stealing?
\end{tcolorbox}

We evaluate five configurations that span the design space of local vs.\ distributed dealing and with vs.\ without stealing:
\begin{description}
    \item[\texttt{GLOBAL\_QUEUE}:] Single shared queue, no per-thread dealing or stealing (baseline).
    \item[\texttt{LOCAL+NoSteal}:] Producer-local dealing (successors stay on the producing thread), no stealing.
    \item[\texttt{LOCAL+STEAL}:] Producer-local dealing with \texttt{CHOOSE\_TWO} stealing when idle.
    \item[\texttt{DISTRIBUTED+NoSteal}:] Round-robin dealing (successors sent to other threads), no stealing.
    \item[\texttt{DISTRIBUTED+STEAL}:] Round-robin dealing with \texttt{CHOOSE\_TWO} stealing when idle.
\end{description}

% % ######% ######% ######% ######% ######% ######% ######% ######% ######% ######
\subsection{Experimental Results}
\label{sec:combined-results}

We evaluate these five configurations on the same four workloads as before, using 50 million tuples (4.8\,GB) with thread counts from 1 to 64.
For thread counts $\leq 32$, processes are pinned to NUMA node~0 via \texttt{numactl}.

\subsubsection{Stateless Workload}

Table~\ref{tab:ds-stateless} shows the throughput.

\begin{table}[h]
\centering
\caption{Stateless workload throughput (MB/s). Bold = best per thread count.}
\label{tab:ds-stateless}
\begin{tabular}{l|rrrrrrr}
\hline
\textbf{Config} & \textbf{1T} & \textbf{2T} & \textbf{4T} & \textbf{8T} & \textbf{16T} & \textbf{32T} & \textbf{64T} \\
\hline
\texttt{GLOBAL\_QUEUE}       & 77 & 149 & 292 & 576 & \textbf{935} & \textbf{771} & \textbf{578} \\
\texttt{LOCAL+NoSteal}       & 76 & \textbf{150} & \textbf{295} & \textbf{581} & 586 & 677 & 531 \\
\texttt{LOCAL+STEAL}         & 76 & 150 & 289 & 571 & 586 & 678 & 534 \\
\texttt{DISTRIBUTED+NoSteal} & 76 & 148 & 287 & 555 & 571 & 670 & 531 \\
\texttt{DISTRIBUTED+STEAL}   & 76 & 148 & 283 & 550 & 465 & 707 & 527 \\
\hline
\end{tabular}
\end{table}

The key observations are as follows.
%
\begin{itemize}
    \item \textbf{\texttt{GLOBAL\_QUEUE} dominates at 16--64T}, peaking at 935\,MB/s at 16T (12.1$\times$).
    Since this workload produces no successor tasks, per-thread queue infrastructure adds overhead without benefit.

    \item \textbf{\texttt{LOCAL+NoSteal} matches or slightly leads at 1--8T} (581 vs.\ 576\,MB/s at 8T).
    At low thread counts, the per-thread queue overhead is negligible and admission round-robin distributes work as well as the global queue.

    \item \textbf{Adding stealing to local dealing has no effect} (\texttt{LOCAL+STEAL} $\approx$ \texttt{LOCAL+NoSteal}, within 2\%).
    With no successor tasks in per-thread queues, there is nothing to steal.

    \item \textbf{Distributed dealing consistently underperforms local dealing} by 3--5\% at 4--8T.
    The round-robin dispatch overhead (atomic counter increment per successor task) adds up even though the successor tasks themselves are rare.

    \item \textbf{\texttt{DISTRIBUTED+STEAL} shows anomalous behavior at 16T} (465\,MB/s, 20\% below \texttt{DISTRIBUTED+NoSteal}).
    The $O(1)$ stealing attempts on empty queues create cache line traffic that interferes with the admission queue contention at this thread count.
\end{itemize}

\subsubsection{High-Frequency Window Workload}

Table~\ref{tab:ds-highfreq} shows the throughput.

\begin{table}[h]
\centering
\caption{High-frequency window workload throughput (MB/s). Bold = best per thread count.}
\label{tab:ds-highfreq}
\begin{tabular}{l|rrrrrrr}
\hline
\textbf{Config} & \textbf{1T} & \textbf{2T} & \textbf{4T} & \textbf{8T} & \textbf{16T} & \textbf{32T} & \textbf{64T} \\
\hline
\texttt{GLOBAL\_QUEUE}       & 39 & 51 & 75 & 103 & 118 & 130 & 40 \\
\texttt{LOCAL+NoSteal}       & 38 & \textbf{53} & \textbf{87} & \textbf{123} & \textbf{129} & 130 & 40 \\
\texttt{LOCAL+STEAL}         & 39 & 48 & 73 & 108 & 108 & \textbf{131} & \textbf{45} \\
\texttt{DISTRIBUTED+NoSteal} & 38 & 52 & 79 & 115 & 120 & 117 & 27 \\
\texttt{DISTRIBUTED+STEAL}   & 39 & 49 & 72 & 100 & 105 & 123 & 54 \\
\hline
\end{tabular}
\end{table}

The key observations are as follows.
%
\begin{itemize}
    \item \textbf{\texttt{LOCAL+NoSteal} leads at 2--16T}, achieving 123\,MB/s at 8T vs.\ 103 for \texttt{GLOBAL\_QUEUE} (+19\%).
    This is the strongest evidence that producer-local dealing improves throughput for workloads with frequent successor tasks ($\sim$5 million window emissions).
    The cache locality benefit of keeping successor tasks on the producing thread is significant.

    \item \textbf{Adding stealing to local dealing hurts at 2--16T} but helps slightly at 32--64T.
    \texttt{LOCAL+STEAL} drops to 108\,MB/s at 8T ($-$12\% vs.\ \texttt{LOCAL+NoSteal}).
    Stealing moves cache-hot window emission tasks to threads with cold caches.
    At 32T, \texttt{LOCAL+STEAL} ties \texttt{LOCAL+NoSteal} (131 vs.\ 130), and at 64T it leads (45 vs.\ 40) because cross-NUMA imbalance benefits from rebalancing.

    \item \textbf{Distributed dealing is consistently worse than local dealing.}
    \texttt{DISTRIBUTED+NoSteal} achieves 115\,MB/s at 8T (vs.\ 123 for \texttt{LOCAL+NoSteal}), confirming that round-robin dispatch of successor tasks destroys the cache locality that the high-frequency workload depends on.

    \item \textbf{At 64T (cross-NUMA), \texttt{DISTRIBUTED+STEAL} leads} with 54\,MB/s.
    When all threads span both NUMA nodes, distributing work and stealing from remote nodes is better than keeping tasks local on one NUMA node.
    However, this comes at the cost of much lower absolute performance (54 vs.\ 130\,MB/s at 32T).
\end{itemize}

\subsubsection{Join Workload}

Table~\ref{tab:ds-join} shows the throughput.

\begin{table}[h]
\centering
\caption{Join workload throughput (MB/s). Bold = best per thread count.}
\label{tab:ds-join}
\begin{tabular}{l|rrrrrrr}
\hline
\textbf{Config} & \textbf{1T} & \textbf{2T} & \textbf{4T} & \textbf{8T} & \textbf{16T} & \textbf{32T} & \textbf{64T} \\
\hline
\texttt{GLOBAL\_QUEUE}       & 31 & \textbf{41} & \textbf{67} & \textbf{103} & 108 & 85 & 60 \\
\texttt{LOCAL+NoSteal}       & 31 & 40 & 65 & 100 & \textbf{113} & 93 & 70 \\
\texttt{LOCAL+STEAL}         & 31 & 41 & 61 & 86 & 99 & \textbf{102} & 64 \\
\texttt{DISTRIBUTED+NoSteal} & 30 & 41 & 62 & 99 & 110 & 93 & \textbf{72} \\
\texttt{DISTRIBUTED+STEAL}   & 31 & 41 & 64 & 99 & 109 & 95 & 71 \\
\hline
\end{tabular}
\end{table}

The key observations are as follows.
%
\begin{itemize}
    \item \textbf{The join workload shows a crossover between strategies as thread count increases.}
    \texttt{GLOBAL\_QUEUE} leads at 4--8T (67/103\,MB/s) where the join's hash table contention dominates.
    \texttt{LOCAL+NoSteal} takes over at 16T (113\,MB/s, +5\%) as the successor tasks from window emissions benefit from cache locality.
    \texttt{LOCAL+STEAL} wins at 32T (102\,MB/s, +20\% over \texttt{GLOBAL\_QUEUE}'s 85\,MB/s) where hyperthreading creates pipeline imbalance that stealing can address.

    \item \textbf{Stealing hurts at 4--16T but helps at 32T.}
    \texttt{LOCAL+STEAL} drops to 86\,MB/s at 8T ($-$14\% vs.\ \texttt{LOCAL+NoSteal}).
    The join's successor tasks carry intermediate results that benefit from staying on the producing thread's cache.
    At 32T, the three pipelines (left source, right source, join) create structural imbalance: some threads finish their source pipeline early and idle, while others are still processing the join.
    Stealing redistributes the backed-up join tasks, improving throughput by 20\%.

    \item \textbf{\texttt{DISTRIBUTED+NoSteal} is competitive at 16--64T} (110/93/72\,MB/s).
    Unlike the high-frequency workload, the join's successor tasks are less cache-sensitive (join inputs are new data from the hash table, not continuation of the same buffer), so distributing them does less harm.

    \item \textbf{At 64T (cross-NUMA), per-thread strategies lead} (\texttt{DISTRIBUTED+NoSteal} at 72, \texttt{LOCAL+NoSteal} at 70, vs.\ \texttt{GLOBAL\_QUEUE} at 60\,MB/s).
    The join's multi-pipeline structure benefits from per-thread queues that keep work on the same NUMA node.
\end{itemize}

\subsubsection{Mixed Concurrent Workload}

Table~\ref{tab:ds-mixed} shows the throughput (computed as total data processed divided by wall-clock time).

\begin{table}[h]
\centering
\caption{Mixed workload throughput (MB/s, computed from wall-clock time). Bold = best per thread count.}
\label{tab:ds-mixed}
\begin{tabular}{l|rrrrrrr}
\hline
\textbf{Config} & \textbf{1T} & \textbf{2T} & \textbf{4T} & \textbf{8T} & \textbf{16T} & \textbf{32T} & \textbf{64T} \\
\hline
\texttt{GLOBAL\_QUEUE}       & 49 & 77 & 130 & 195 & 227 & \textbf{240} & 83 \\
\texttt{LOCAL+NoSteal}       & 49 & \textbf{77} & \textbf{136} & \textbf{212} & \textbf{233} & 239 & 86 \\
\texttt{LOCAL+STEAL}         & 49 & 76 & 128 & 191 & 200 & 238 & 96 \\
\texttt{DISTRIBUTED+NoSteal} & 49 & 76 & 128 & 197 & 220 & 214 & 55 \\
\texttt{DISTRIBUTED+STEAL}   & 48 & 76 & 123 & 178 & 188 & 217 & \textbf{112} \\
\hline
\end{tabular}
\end{table}

The key observations are as follows.
%
\begin{itemize}
    \item \textbf{\texttt{LOCAL+NoSteal} leads at 2--16T}, peaking at 233\,MB/s at 16T vs.\ 227 for \texttt{GLOBAL\_QUEUE} (+3\%).
    At 8T the gap is larger: 212 vs.\ 195\,MB/s (+9\%).
    The mixed workload combines a lightweight filter (no successors) with a heavy aggregation ($\sim$5M successors), and keeping the aggregation's successor tasks local avoids contention with the filter query's workers.

    \item \textbf{Stealing on top of local dealing hurts consistently at 4--16T.}
    \texttt{LOCAL+STEAL} achieves 191\,MB/s at 8T ($-$10\% vs.\ \texttt{LOCAL+NoSteal}).
    The filter query's idle workers steal aggregation successor tasks, processing them with cold caches.
    This is the same pattern observed in the high-frequency workload: stealing cache-hot window emissions to cache-cold threads degrades performance.

    \item \textbf{Distributed dealing is worse than local dealing at all single-NUMA thread counts} (4--32T).
    \texttt{DISTRIBUTED+NoSteal} achieves 197\,MB/s at 8T (vs.\ 212 for \texttt{LOCAL+NoSteal}), confirming that round-robin dispatch of successor tasks is suboptimal for workloads that benefit from cache locality.

    \item \textbf{At 64T (cross-NUMA), \texttt{DISTRIBUTED+STEAL} leads} with 112\,MB/s vs.\ 83--96 for all other configurations.
    This is the only scenario where the combination of distributing work and stealing provides a net benefit.
    With 64 threads spanning two NUMA nodes, the workload imbalance between the fast filter and slow aggregation is amplified by non-uniform memory access, and aggressive rebalancing compensates.
\end{itemize}

\subsection{Summary}

\begin{itemize}
    \item \textbf{\texttt{LOCAL+NoSteal} is the best configuration for successor-heavy workloads at 2--16T.}
    It leads on the high-frequency workload (+19\% at 8T), the mixed workload (+9\% at 8T), and the join workload (+5\% at 16T).
    Producer-local dealing ensures that successor tasks are processed with warm L1/L2 caches, while the absence of stealing preserves this locality.

    \item \textbf{\texttt{GLOBAL\_QUEUE} remains the best for workloads without successor tasks.}
    On the stateless workload, it peaks at 935\,MB/s at 16T---60\% higher than any per-thread configuration.
    Its simplicity and perfect load balancing make it the optimal choice when the workload is dominated by admission tasks.

    \item \textbf{Stealing helps only at $\geq 32$T on multi-pipeline workloads.}
    \texttt{LOCAL+STEAL} achieves 102\,MB/s at 32T on the join workload (+20\% over \texttt{GLOBAL\_QUEUE}), because the three-pipeline structure creates structural imbalance that dealing alone cannot resolve.
    At lower thread counts, stealing consistently hurts by breaking cache locality.

    \item \textbf{Distributed dealing is always inferior to local dealing.}
    Across all four workloads and all thread counts, \texttt{DISTRIBUTED+NoSteal} trails \texttt{LOCAL+NoSteal}.
    The round-robin dispatch overhead and loss of cache locality for successor tasks are never compensated by the marginally better load distribution.
    The sole exception is 64T cross-NUMA, where \texttt{DISTRIBUTED+STEAL} leads on the mixed workload.

    \item \textbf{The cost of cache-cold stealing is workload-dependent.}
    On the join workload, stealing reduces throughput by 14\% at 8T (successor tasks carry join hash table state).
    On the high-frequency workload, the reduction is 12\% at 8T (successor tasks carry window aggregation state).
    On the mixed workload, it is 10\% at 8T.
    The common pattern: the more compute-intensive the successor task, the higher the penalty from processing it on a cache-cold thread.

    \item \textbf{The optimal configuration depends on thread count and workload:}
    \begin{itemize}
        \item 1--8T: \texttt{LOCAL+NoSteal} (or \texttt{GLOBAL\_QUEUE} if no successors)
        \item 16T: \texttt{LOCAL+NoSteal} for successor-heavy, \texttt{GLOBAL\_QUEUE} for stateless
        \item 32T: \texttt{LOCAL+STEAL} for multi-pipeline, \texttt{GLOBAL\_QUEUE} for stateless
        \item 64T (cross-NUMA): \texttt{DISTRIBUTED+STEAL} for mixed, \texttt{GLOBAL\_QUEUE} for stateless
    \end{itemize}
\end{itemize}
